\scanpage{128}\mspage{73}
Il s'ensuit notamment que \uline{les foncteurs} $i,j$ \uline{sont pleinement fidèles}. \deleted{Pour que l'\unsure{élém$^{\mathrm t}$} $\xi$ de $\mathrm{HOT}_W(S)$, défini par l'objet $X$ de $\mathrm{Cat}_{/S}$ soit dans l'image essentielle de $i$, il \unsure{f. et s.} \unsure{que}}

\textcircled{V} \deleted{\unsure{Des}}On a des foncteurs évidents
\[
 \mathrm{HOT}_W(S)\unsure{\longrightarrow|\text{---}}\text{\uline{Hom}}\,(S^\circ,\mathrm{Hot}_W)
 \tag*{$(*)$}
\]
\[
 \mathrm{HOT}^\circ_W(S)\longrightarrow\text{\uline{Hom}}\,(S,\mathrm{Hot}_W),
 \tag*{$(*')$}
\]
le premier s'interprète sur les \deleted{trois}\inserted{quatre} catégories $(*)$ \inserted{indifféremment} \deleted{$\text{\uline{Fibsc}}\,S$} (bordantes gauches) de \unsure{page|p.}~70, \deleted{\unsure{qui}} servant à définir $\mathrm{HOT}_W(S)$, par
\[
 s\longrightarrow\mathrm{hot}_W(X_s)
\]
dans le cas des trois catégories distinctes de $\mathrm{Cat}_{/S}$, et par
\[
 s\longmapsto\mathrm{hot}_W({}_{s\backslash}X)
\]
dans le cas de $\mathrm{Cat}_{/S}$. La dépendance fonctorielle par rapport à $s$ est évidente dans le cas de $\mathrm{Cat}_{/S}$ (donné par $s\mapsto\mathrm{hot}_W({}_{s\backslash}X)$), et dans les cas de $\text{\uline{Fibsc}}\,S$, $\text{\uline{Fib}}\,S$, où on dispose, pour $u:s\unsure{\text{---}|\longrightarrow}t$ dans $S$, du changement de base
\[
 u^*:X_{\unsure{b|t}}\text{---}X_s.
\]
Dans le cas de $\text{\uline{Liss}}\,S$, la dépendance fonctorielle se déduit de celle \deleted{\unsure{du}} de $\mathrm{hot}_W({}_{s\backslash}X)$, grâce au fait que l'inclusion
\[
 X_s\longrightarrow{}_{s\backslash}X
\]
est alors une $W$-équivalence.

On a une interprétation duale de l'application
% archive sheet number: 127
